% !TeX root = ../papers/nc-p3.tex
\section{三因子対称境界と有限切断代数}
\label{sec:truncated-quadric-boundary}
三次商の双対基底から作る前節のテンソルは、有限切断代数のtrivial extensionとして理解できる。
この記述は新しいAS正則性の十分条件ではないが、個々の大行列に依存せず境界問題を整理する。

\begin{lemma}[四層二次代数のKoszul性]\label{lem:four-layer-koszul}
頂点$0,1,2,3$を持ち、矢が$i$から$i+1$へだけ向かう箙の二次商代数はKoszulであり、
大域次元は3以下である。
\end{lemma}
\begin{proof}
経路の長さは始点と終点の差に等しく、最大3である。
単純加群の極小射影分解では、$p$番目の項の内部次数は$p$以上であり、
始点から終点へ向かう正次数成分の積として長さ3以下でなければならない。
従って$p>3$の項は零である。
第1項は次数1生成、第2項は二次関係による次数2生成である。
第3項は内部次数3以外を持ち得ないので、分解は線型である。
\end{proof}

$L\subset\operatorname{Sym}^2(V^*)$を四次元とし、
$B=\operatorname{Sym}(V^*)/(L)$が$\dim B_1=4$, $\dim B_2=6$, $\dim B_3=4$を満たすとする。
三次双対基底$b_0,\ldots,b_3$から、因子$0,1,3$が対称なテンソル
$w_{abjc}=(b_j)_{abc}$を作る。
有限有向代数$A_L$を、順序づけた頂点$3,0,1,2$に対して
\[
 e_iA_Le_j=B_{\operatorname{pos}(j)-\operatorname{pos}(i)}
 \quad(\operatorname{pos}(i)\leq\operatorname{pos}(j))
\]
とし、逆向き成分は零、積は$B$の積で定める。
ここで$\operatorname{pos}(3)=0$, $\operatorname{pos}(0)=1$,
$\operatorname{pos}(1)=2$, $\operatorname{pos}(2)=3$である。
$A_L$は四層二次代数なのでKoszulであり、$\dim A_L=32$である。
$D A_L=\operatorname{Hom}_k(A_L,k)$の$B_r$双対成分に次数$4-r$を与え、
通常のtrivial extension
\[
 T_L=A_L\ltimes D A_L,
 \qquad (a,f)(b,g)=(ab,a\cdot g+f\cdot b)
\]
にこの次数を入れる。

\begin{proposition}[三因子対称Frobenius候補のtrivial extension表示]
\label{prop:partial-symmetric-trivial-extension}
上記$w$がテンソル復元定理の有限Frobenius条件を満たすなら、
復元した二次双対を四頂点で折り畳んだ代数$E(w)$は、次数付き代数として$T_L$と同型である。
従ってこの部分のKoszul判定は、$T_L$のKoszul判定と一致する。
この主張に$B$の完全交叉性は仮定しない。
\end{proposition}
\begin{proof}
$T_L$の次数一成分は、三本の前向き辺上の$B_1$と戻り辺$2\to3$上の$B_3^*$である。
trivial extensionのtraceは、$D A_L$成分を$1\in A_L$で評価する写像であり、
相補次数の積を完全に対合させる。
三つの$B_1$と一つの$B_3^*$の積のtraceは
\[
 (f,a,b,c)\longmapsto f(abc)
\]
である。選んだ基底でのこの四階テンソルが、まさに$w$とその巡回回転である。
各始点の$T_L$の次数別次元は$(1,4,6,4,1)$となる。

$w$の一因子階数4と隣接二対二階数6により、次数一成分の積は
$T_L$の次数2,3,4を全て生成する。実際、その積写像の階数は、完全trace対合を使うと
対応する一因子・二因子縮約の階数に等しい。
従って$T_L$は次数一生成であり、その二次関係は$R_i(w)^\perp$そのものである。
よって二次表示から全射$E(w)\twoheadrightarrow T_L$を得る。
有限Frobenius仮定から両者の全次元は64で等しいので、これは同型である。
\end{proof}

\begin{remark}[既存のhigher zigzag構成との関係]
Grant\cite[Definitions 2.3, 2.5]{GrantZigzag}は、有限大域次元のKoszul代数$\Lambda$に対し、
その二次双対の符号捩れtrivial extensionをhigher zigzag algebraとして定義している。
ここでは$\Lambda=A_L^!$と取れる。
四層の長さ符号は頂点上の$\operatorname{diag}(1,-1,1,-1)$による内部自己同型なので、
その符号捩れは次数零の可逆元で除ける。
従って上の$T_L$は、この既知構成の四層の場合としても記述できる。
trivial extensionという構成自体の新規性や、それだけからのKoszul性は主張しない。
本研究で必要なのは、基点境界で二次表示とKoszul性がどの段階で破れるかの判定である。
\end{remark}

\begin{lemma}[一般点の有限次数障害による族全体の排除]
\label{lem:generic-quintic-exclusion}
既約なパラメータ空間$X$上の正則な射影テンソル族で、
全ての隣接二対二階数が6以下であるものを取り、
有限Frobenius条件と全四次Koszul検査を通る部分$U\subset X$が空でないとする。
$U$は開部分である。
一般点で五次Koszul微分の階数が
$\operatorname{rk}d_2=a$, $\operatorname{rk}d_3=b$なら、
全ての$U$上の点で$\dim H_2\geq120-a-b$である。
特に$a+b<120$なら、この族は標準型Koszul AS対象を一つも含まない。
\end{lemma}
\begin{proof}
隣接階数の上界により、有限Frobenius条件は$X$内で開であり、四次検査の開性と合わせて$U$は開である。
$U$上では必要な次数2,3,4の商空間は一定階数であり、
$K_2,K_3,K_1$はそれぞれ階数120,40,140のベクトル束になる。
二つの微分はベクトル束の射なので、その階数は一般点の値を超えない。
従って
$\dim H_2=120-\operatorname{rk}d_2-\operatorname{rk}d_3\geq120-a-b$である。
標準型Koszul AS点は必ず$U$に入るので、最後の主張も従う。
\end{proof}

\begin{theorem}[係数を動かした基点境界族の一様なAS排除]
\label{thm:generic-one-point-uniform-exclusion}
$c_0c_1c_2\ne0$とし、
\[
 L_{c,t}=\langle e_i^2-e_ie_3\ (i=0,1,2),\
                    e_3(c_0e_0+c_1e_1+c_2e_2+t e_3)\rangle
\]
から、四座標$012,133,233,333$を正規化した三次双対基底によるテンソル$w_{c,t}$を作る。
$t=0$の係数族は、全て$Z_{\mathrm{sym}}$に属するが、標準型Koszul AS対象を一つも含まない。
より正確には、有限Frobenius条件と全四次検査を通る任意の係数値で、
始点0の五次Koszulホモロジーの次元は7以上である。

一般点では有限Frobenius条件と全四次検査が成立し、
五次の$H_2$次元は$(7,0,0,0)$、$H_3$は全始点で零である。
始点0の七次元空間のうち六次元は、前節と同じ平方型関係と左乗法の核から来る。
ただし係数全域の有限Frobenius層や、一般の四二次式の基点境界を分類したとは主張しない。
\end{theorem}
\begin{proof}
最後の式を$c_0$で割り、$u=c_1/c_0$, $v=c_2/c_0$と正規化する。
$t=0$のテンソル族は$X=(\mathbb G_m)^2$上の正則族になる。
三次の16本の消去式と四本の正規化式を合わせた行列の行列式は$-2uv$である。
従って指定三次基底は$X$全体で定まり、特にその第3ベクトルは$333$のcubeである。
他の$t$も含む行列式は$-2c_0^2c_1c_2$なので、$t$方向にも同じフレームが定まる。
各$c$について、最初の三二次式の八点上で最後の式が消える$t$は有限個である。
それ以外の$t$では基点自由完全交叉になるから、$w_{c,0}\in Z_{\mathrm{sym}}$である。

三因子が対称であり、その二因子に$L_{c,0}$が作用して零になるので、
全ての二対二階数は6以下である。
$u=v=1$では前節の例に一致し、有限Frobenius条件と全四次検査を通る。
従ってそれらを通る部分$U\subset X$は空でない開部分であり、一般点を含む。
同じ標本では、始点1,2,3の五次複体も正確である。
その正確性も一定次元の複体の開条件なので、一般点では保たれる。

始点0は、二つの不定元を持つ体$\mathbb Q(u,v)$上で計算する。
復元した六関係を挿入した行列の直交核から商写像と自由単項式sectionを取ると、
$C_{2,5}$は20次元、$C_{1,5}$は35次元であり、
\[
 \operatorname{rk}_{\mathbb Q(u,v)}d_2=77,\qquad
 \operatorname{rk}_{\mathbb Q(u,v)}d_3=36,\qquad d_2d_3=0
\]
を得る。これは有限個の有理数標本からの推論ではない。
二変数多項式の正確なGCDと割り算による有理関数消去の再現は
\path{computations/verify_generic_one_point_boundary.py}に置く。
各算術演算は分母を約分し、階数消去は非零な有理関数だけをpivotに使う。
従って一般点の$H_2$は$120-77-36=7$次元である。

さらに左乗法$m_{x_{1,3}}:C_{2,5}\to C_{1,5}$の一般階数は14で、その核は6次元である。
その核から平方型関係$x_{0,3}\otimes x_{1,3}$を用いて作ったcycle列を$Z$とすると、
$d_2Z=0$、$\operatorname{rk}(d_3\mid Z)=42$が同じ有理関数体上で成立する。
よって一般点でも六つの独立な平方型障害類があり、その商は一次元である。
一般の混合第七類の短い閉じた式は、ここでは与えない。

最後に補題\ref{lem:generic-quintic-exclusion}を適用すれば、全ての$U$上で
$\dim H_2\geq120-77-36=7$となる。
$X\setminus U$では有限Frobenius条件か四次検査を既に失っている。
従って例外的な係数値にも標準型Koszul AS点は存在しない。
全ての非空部分集合$I\subset\{0,1,2\}$について$\sum_{i\in I}c_i\ne0$とすれば
共通基点は$[0:0:0:1]$だけになるが、一様排除にはこの追加条件を必要としない。
\end{proof}
